% ============================================================================
% EXEMPLE - Comment ajouter des PDFs au document LaTeX
% ============================================================================
% 
% Ce fichier montre les différentes façons d'intégrer vos PDFs
% dans le document main.tex
%
% ============================================================================

% --- OPTION 1 : Inclure un PDF complet (toutes les pages) ---

\subsection{Rapport d'évaluation complet}
\includepdf[pages=-]{pdfs/evaluation_report.pdf}


% --- OPTION 2 : Inclure un sous-ensemble de pages ---

\subsection{Architecture du modèle (pages 1-3)}
\includepdf[pages=1-3]{pdfs/model_architecture.pdf}


% --- OPTION 3 : Inclure des pages spécifiques ---

\subsection{Résultats clés}
\includepdf[pages={1,5,10}]{pdfs/results_summary.pdf}


% --- OPTION 4 : Inclure une seule page ---

\subsection{Graphique principal}
\includepdf[pages=1]{pdfs/main_chart.pdf}


% --- OPTION 5 : Avec texte descriptif avant le PDF ---

\subsection{Matrices de confusion}

Les matrices de confusion montrent la performance détaillée du modèle Run B
sur les données de test. Elles révèlent :
\begin{itemize}
    \item La classe \textit{destroyed} est bien identifiée
    \item La classe \textit{no-damage} atteint 99\% de précision
    \item La classe \textit{minor-damage} reste le point faible
\end{itemize}

\includepdf[pages=-]{pdfs/confusion_matrices.pdf}


% --- OPTION 6 : Plusieurs PDFs en séquence ---

\subsection{Visualisations Grad-CAM}

\subsubsection{Exemples : bâtiments non-endommagés}
\includepdf[pages=1-5]{pdfs/grad_cam_part1.pdf}

\subsubsection{Exemples : bâtiments endommagés}
\includepdf[pages=1-5]{pdfs/grad_cam_part2.pdf}

\subsubsection{Analyse des patterns visuels}
\includepdf[pages=-]{pdfs/grad_cam_analysis.pdf}


% --- OPTION 7 : Avec régulation des espacements ---

\subsection{Données brutes}
\newpage
\includepdf[pages=-]{pdfs/raw_data.pdf}


% --- OPTION 8 : Avec titre personnalisé ---

\subsection{Validation expérimentale}

Le document ci-dessous valide la reproductibilité de la data foundation :

\includepdf[pages=-]{pdfs/data_validation.pdf}

La validation confirme que :
\begin{itemize}
    \item Tous les splits sont sans chevauchement
    \item Les métadonnées sont complètes
    \item Les chemins de fichiers sont corrects
\end{itemize}


% ============================================================================
% COMMENT UTILISER CET EXEMPLE
% ============================================================================
%
% 1. Copiez les sections qui vous conviennent
% 2. Remplacez les noms de fichiers (eg. "evaluation_report.pdf")
%    par vos vrais noms de PDFs
% 3. Collez-les dans le fichier main.tex, dans la section \appendix
%
% Exemple complet pour main.tex :
%
% \appendix
% 
% \section{Annexes - Résultats}
% 
% \subsection{Rapport d'évaluation}
% \includepdf[pages=-]{pdfs/evaluation_report.pdf}
% 
% \subsection{Grad-CAM}
% \includepdf[pages=-]{pdfs/grad_cam_results.pdf}
%
% ============================================================================
